\chapter[Le Condizioni e i Bivi nel City Builder Marziano]{Le Condizioni e i Bivi\\ nel City Builder Marziano}
\label{chap:condizioni}

\begin{quote}
\itshape
\small
``Un gestionale diventa intelligente nel momento in cui impara a guardare il magazzino prima di spendere.''
\end{quote}
\vspace{1em}

\noindent\textbf{Riferimento Lezione:} 5 Ottobre 2026 -- Terza Lezione di Laboratorio\\
\textbf{Progetto Guida Annuale:} \textbf{\color{PiazzaPrimary}Mars City Builder (Lite)}\\
\textbf{Competenze chiave:} Formulazione di condizioni booleane, uso degli operatori relazionali, struttura di selezione \texttt{SE \dots\ ALLORA \dots\ ALTRIMENTI}, disegno di Flow Chart con le forme standard.

\vspace{1.5em}

% ------------------------------------------------------------------------------
\section{Perché un Gestionale ha Bisogno di Bivi?}
% ------------------------------------------------------------------------------

Con la \textbf{struttura sequenziale} (Capitolo~\ref{chap:sequenza}) le istruzioni vengono eseguite una dietro l'altra, senza mai deviare. Immaginiamo però di programmare il nostro \textbf{city builder su Marte} usando solo la sequenza.

Il giocatore clicca sul pulsante \emph{``Costruisci Cupola Idroponica''}. Se il programma non controlla prima il magazzino, piazza l'edificio sulla mappa anche se il giocatore ha 0 unità di metallo: il bilancio scende a $-50$ e la partita si rompe.

\begin{concetto}[La Selezione]
Un programma diventa capace di reagire alla situazione quando sa \textbf{osservare lo stato delle risorse prima di agire}. Questa struttura logica si chiama \textbf{selezione} (o bivio condizionale): in base a una condizione, il programma sceglie quale strada percorrere.
\end{concetto}

% ------------------------------------------------------------------------------
\section{Cos'è una Condizione? La Logica Booleana}
% ------------------------------------------------------------------------------

Una \textbf{condizione} è un'espressione di confronto che ammette solo \textbf{due} risultati possibili: \textbf{VERO} oppure \textbf{FALSO}. Non esistono vie di mezzo: o hai abbastanza metallo o non ce l'hai. Questa logica a due valori si chiama \textbf{logica booleana}, dal matematico George Boole.

\begin{esempio}[Condizioni nel City Builder]
\begin{itemize}
  \item \texttt{metallo >= 50}: con metallo $= 70$ è \textbf{VERO}; con metallo $= 30$ è \textbf{FALSO}.
  \item \texttt{energia\_prodotta >= energia\_richiesta}: con $100 \geq 80$ è \textbf{VERO}; con $60 \geq 80$ è \textbf{FALSO}.
  \item \texttt{coloni == 0}: se i coloni sono 0 è \textbf{VERO} (colonia deserta); se sono 12 è \textbf{FALSO}.
\end{itemize}
\end{esempio}

\begin{attenzione}[Frasi che NON Sono Condizioni]
\begin{itemize}
  \item \emph{``La colonia è bella?''}: è soggettiva, il computer non ha gusti.
  \item \emph{``Costruisci qualche casa''}: è vaga, quante case esattamente?
\end{itemize}
\end{attenzione}

% ------------------------------------------------------------------------------
\section{Gli Operatori Relazionali}
% ------------------------------------------------------------------------------

Per confrontare risorse e parametri della colonia usiamo sei operatori relazionali. Nella tabella supponiamo che il giocatore abbia \texttt{crediti = 50}.

\begin{center}
\renewcommand{\arraystretch}{1.35}
\begin{tabular}{|c|c|l|l|c|}
\hline
\rowcolor{PiazzaPrimary!15}
\textbf{Simbolo} & \textbf{Matematica} & \textbf{Significato} & \textbf{Esempio} & \textbf{Risultato} \\
\hline
\texttt{>}  & $>$    & strettamente maggiore & \texttt{crediti > 20}  & \textbf{VERO} \\
\texttt{<}  & $<$    & strettamente minore   & \texttt{crediti < 10}  & \textbf{FALSO} \\
\texttt{>=} & $\geq$ & maggiore o uguale     & \texttt{crediti >= 50} & \textbf{VERO} \\
\texttt{<=} & $\leq$ & minore o uguale       & \texttt{crediti <= 40} & \textbf{FALSO} \\
\texttt{==} & $=$    & uguale a              & \texttt{crediti == 50} & \textbf{VERO} \\
\texttt{!=} & $\neq$ & diverso da            & \texttt{crediti != 50} & \textbf{FALSO} \\
\hline
\end{tabular}
\end{center}

\begin{attenzione}[La Trappola Numero 1: \texttt{=} contro \texttt{==}]
\begin{itemize}
  \item \texttt{=} (\textbf{assegnazione}) è un ordine: prende un valore e lo memorizza nella variabile. \texttt{crediti = 100} deposita 100 crediti nel conto.
  \item \texttt{==} (\textbf{confronto}) è una domanda: chiede se i due valori sono identici. \texttt{crediti == 100} restituisce VERO o FALSO senza toccare la memoria.
\end{itemize}
\end{attenzione}

% ------------------------------------------------------------------------------
\section{La Struttura SE \dots\ ALLORA \dots\ ALTRIMENTI}
% ------------------------------------------------------------------------------

La selezione divide il cammino del programma in due rami \textbf{mutuamente esclusivi}: si esegue il primo oppure il secondo, mai entrambi nello stesso turno. Alla fine i due rami si ricongiungono e il programma prosegue.

\begin{lstlisting}
SE <condizione> ALLORA
    // ramo VERO
    istruzione_A1
    istruzione_A2
ALTRIMENTI
    // ramo FALSO
    istruzione_B1
FINE SE
// il programma si ricongiunge qui
\end{lstlisting}

\begin{consiglio}[La Selezione a Una Via]
Se nel caso FALSO non serve fare niente di speciale, il ramo \texttt{ALTRIMENTI} si può omettere. Esempio: un bonus di fine turno facoltativo.
\begin{lstlisting}
SE colonia_felice == VERO ALLORA
    crediti = crediti + 20
FINE SE
salva_partita()
\end{lstlisting}
\end{consiglio}

\subsection{Il Pattern del City Builder: Verifica e Spesa}

Ogni volta che il giocatore ordina una costruzione, l'algoritmo segue sempre la stessa sequenza logica:
\begin{enumerate}
  \item \textbf{Verifica risorse:} le scorte coprono il costo? (\texttt{SE risorsa >= costo})
  \item \textbf{Aggiornamento magazzino (se VERO):} piazza l'edificio e scala il costo (\texttt{risorsa = risorsa - costo}).
  \item \textbf{Feedback a schermo:} mostra un messaggio chiaro, di conferma o di risorse insufficienti.
\end{enumerate}

% ------------------------------------------------------------------------------
\section{Le Forme del Flow Chart}
% ------------------------------------------------------------------------------

Il \textbf{Flow Chart} (diagramma a blocchi) disegna la logica di un algoritmo. Ogni forma geometrica ha un significato standard:

\begin{center}
\renewcommand{\arraystretch}{1.6}
\begin{tabular}{|>{\centering\arraybackslash}m{3.4cm}|m{4.6cm}|m{4.6cm}|}
\hline
\rowcolor{PiazzaPrimary!15}
\textbf{Forma} & \textbf{Funzione} & \textbf{Esempio} \\
\hline
\tikz\node[fc-term, minimum width=18mm, minimum height=6mm, font=\scriptsize\bfseries]{INIZIO};
  & \textbf{Inizio e fine} del flusso & \texttt{INIZIO} / \texttt{FINE} \\
\hline
\tikz\node[fc-io, minimum height=6mm, font=\scriptsize]{LEGGI / SCRIVI};
  & \textbf{Input e output}: lettura dei dati e stampa a video & \texttt{LEGGI metallo} \\
\hline
\tikz\node[fc-op, minimum height=6mm, font=\scriptsize]{calcolo};
  & \textbf{Operazioni e calcoli}: assegnazioni e modifiche della memoria & \texttt{metallo = metallo - 50} \\
\hline
\tikz\node[fc-dec, font=\scriptsize]{condizione?};
  & \textbf{Condizione e bivio}: due uscite obbligatorie, SÌ e NO & \texttt{metallo >= 50 ?} \\
\hline
\end{tabular}
\end{center}

\begin{concetto}[Le 3 Regole d'Oro del Flow Chart]
\begin{enumerate}
  \item \textbf{Frecce direzionali:} collegano i blocchi dall'alto verso il basso e non restano mai aperte.
  \item \textbf{Etichette sul rombo:} ogni rombo ha due frecce uscenti, etichettate \textbf{SÌ} (VERO) e \textbf{NO} (FALSO).
  \item \textbf{Punto di ricongiunzione:} i due rami si riuniscono prima dell'output finale e di \texttt{FINE}.
\end{enumerate}
\end{concetto}

% ------------------------------------------------------------------------------
\section{Esempio Svolto: la Cupola Idroponica}
% ------------------------------------------------------------------------------

\begin{esempio}[Il Testo del Problema]
Ogni Cupola Idroponica costa \textbf{50 unità di metallo}. L'algoritmo legge il \texttt{metallo\_disponibile} nel magazzino centrale. Se il metallo è \textbf{pari o superiore a 50}, scala 50 unità dal magazzino e imposta il messaggio \emph{``Serra costruita con successo!''}; altrimenti non costruisce nulla e imposta il messaggio \emph{``Metallo insufficiente per la Serra!''}. Alla fine mostra il messaggio e il metallo rimasto.
\end{esempio}

\subsection{Passo 1: Lo Schema I-P-O}

\begin{itemize}
  \item \textbf{Input:} \texttt{metallo\_disponibile} (numero intero $\geq 0$).
  \item \textbf{Elaborazione:} se \texttt{metallo\_disponibile >= 50} scala il costo e prepara il messaggio di successo; altrimenti prepara il messaggio di errore.
  \item \textbf{Output:} il messaggio e il saldo aggiornato del metallo.
\end{itemize}

\subsection{Passo 2: Il Flow Chart}

\begin{center}
\begin{tikzpicture}
  \node[fc-term] (start) at (0,0) {INIZIO};
  \node[fc-io]   (in)    at (0,-1.3) {LEGGI metallo};
  \node[fc-dec]  (dec)   at (0,-2.9) {metallo $\geq$ 50 ?};
  \node[fc-op]   (y1)    at (-3.4,-4.6) {metallo = metallo $-$ 50};
  \node[fc-op]   (y2)    at (-3.4,-5.9) {msg = ``Serra costruita!''};
  \node[fc-op]   (n1)    at (3.4,-4.6) {msg = ``Metallo insufficiente!''};
  \node[fc-join] (join)  at (0,-7.1) {};
  \node[fc-io]   (out)   at (0,-8.2) {SCRIVI msg, metallo};
  \node[fc-term] (end)   at (0,-9.5) {FINE};
  \draw[fc-arrow] (start) -- (in);
  \draw[fc-arrow] (in) -- (dec);
  \draw[fc-arrow] (dec.west) -| node[fc-label, pos=0.25, above] {SÌ} (y1.north);
  \draw[fc-arrow] (dec.east) -| node[fc-label, pos=0.25, above] {NO} (n1.north);
  \draw[fc-arrow] (y1) -- (y2);
  \draw[fc-arrow] (y2.south) |- (join);
  \draw[fc-arrow] (n1.south) |- (join);
  \draw[fc-arrow] (join) -- (out);
  \draw[fc-arrow] (out) -- (end);
\end{tikzpicture}
\end{center}

\subsection{Passo 3: Lo Pseudocodice}

\begin{lstlisting}
// 1. Acquisizione del dato dal magazzino
LEGGI metallo_disponibile

// 2. Controllo di fattibilita' e spesa
SE metallo_disponibile >= 50 ALLORA
    metallo_disponibile = metallo_disponibile - 50
    messaggio = "Serra costruita con successo!"
ALTRIMENTI
    messaggio = "Metallo insufficiente per la Serra!"
FINE SE

// 3. Notifica al giocatore
SCRIVI messaggio
SCRIVI "Metallo rimasto in magazzino: ", metallo_disponibile
\end{lstlisting}

% ------------------------------------------------------------------------------
\newpage
\section{Esercizi di Laboratorio: Disegna il Flow Chart}
% ------------------------------------------------------------------------------

Per ciascun esercizio disegna sul quaderno, con penna e righello, il \textbf{Flow Chart completo} rispettando le forme standard.

\begin{esercizio}[Esercizio 1: Il Bilancio Energetico della Colonia]
\emph{Il sole rosso di Marte sta calando dietro i crateri di Tharsis. I coloni rientrano nelle cupole, accendono i riscaldatori e i purificatori d'aria passano alla modalità notturna. Come responsabile dell'energia, devi verificare se la carica dei pannelli solari basterà per tutta la notte.}

\begin{itemize}
  \item Se l'energia prodotta copre o supera il fabbisogno della base, la rete è stabile e la colonia passa la notte in sicurezza.
  \item Se la produzione è insufficiente, scatta l'allarme blackout: bisogna staccare macchinari e fabbriche.
  \item Il programma mostra a schermo lo stato della rete.
\end{itemize}
\end{esercizio}

\begin{esercizio}[Esercizio 2: Nuovi Coloni e Razioni di Cibo]
\emph{I retrorazzi della navetta ``Ares Shuttle'' sollevano una nube di polvere rossa: sono sbarcati nuovi coloni. Prima di aprire i portelloni, l'amministratore consulta i silos, sapendo che ogni colono consuma esattamente una cassa di cibo a turno.}

\begin{itemize}
  \item Se le casse bastano (o avanzano) per tutti i coloni, il cibo viene distribuito e le scorte diminuiscono della quantità consumata.
  \item Se le scorte non bastano per tutti, la distribuzione non avviene e scatta l'allarme carestia.
  \item Il programma mostra l'esito e il saldo finale delle scorte.
\end{itemize}
\end{esercizio}

\begin{esercizio}[Esercizio 3: L'Indice di Felicità della Colonia]
\emph{È venerdì sera nella cupola centrale e si tiene l'assemblea cittadina. I sensori sociali registrano l'indice di gradimento della popolazione, da 0 a 100.}

\begin{enumerate}
  \item \textbf{Colonia fiorente:} se il gradimento raggiunge o supera 75, la colonia incassa un bonus di 50 crediti.
  \item \textbf{Rivolta:} se il gradimento scende sotto 40, scatta lo sciopero e la produzione si blocca.
  \item \textbf{Colonia stabile:} in tutti i casi intermedi (da 40 a 74) la vita prosegue regolarmente.
\end{enumerate}
Il programma comunica lo stato finale. \textbf{Suggerimento:} usa due rombi a cascata, il secondo sul ramo NO del primo.
\end{esercizio}

% ------------------------------------------------------------------------------
\newpage
\section{Soluzioni Commentate}
% ------------------------------------------------------------------------------

\begin{soluzione}[Esercizio 1: Il Bilancio Energetico]
\begin{center}
\begin{tikzpicture}[font=\small]
  \node[fc-term] (start) at (0,0) {INIZIO};
  \node[fc-io]   (in)    at (0,-1.3) {LEGGI prodotta, richiesta};
  \node[fc-dec]  (dec)   at (0,-3.0) {prodotta $\geq$ richiesta ?};
  \node[fc-op]   (y1)    at (-3.6,-4.8) {stato = ``RETE STABILE''};
  \node[fc-op]   (n1)    at (3.6,-4.8) {stato = ``BLACKOUT: \\ stacca le fabbriche!''};
  \node[fc-join] (join)  at (0,-6.2) {};
  \node[fc-io]   (out)   at (0,-7.3) {SCRIVI stato};
  \node[fc-term] (end)   at (0,-8.6) {FINE};
  \draw[fc-arrow] (start) -- (in);
  \draw[fc-arrow] (in) -- (dec);
  \draw[fc-arrow] (dec.west) -| node[fc-label, pos=0.25, above] {SÌ} (y1.north);
  \draw[fc-arrow] (dec.east) -| node[fc-label, pos=0.25, above] {NO} (n1.north);
  \draw[fc-arrow] (y1.south) |- (join);
  \draw[fc-arrow] (n1.south) |- (join);
  \draw[fc-arrow] (join) -- (out);
  \draw[fc-arrow] (out) -- (end);
\end{tikzpicture}
\end{center}
``Copre o supera'' diventa $\geq$: con produzione esattamente uguale al fabbisogno la rete è stabile.
\end{soluzione}

\begin{soluzione}[Esercizio 2: Nuovi Coloni e Razioni]
\begin{center}
\begin{tikzpicture}[font=\small]
  \node[fc-term] (start) at (0,0) {INIZIO};
  \node[fc-io]   (in)    at (0,-1.3) {LEGGI scorte, coloni};
  \node[fc-dec]  (dec)   at (0,-2.9) {scorte $\geq$ coloni ?};
  \node[fc-op]   (y1)    at (-3.6,-4.6) {scorte = scorte $-$ coloni};
  \node[fc-op]   (y2)    at (-3.6,-5.9) {esito = ``Coloni sazi!''};
  \node[fc-op]   (n1)    at (3.6,-4.6) {esito = ``CARESTIA!''};
  \node[fc-join] (join)  at (0,-7.1) {};
  \node[fc-io]   (out)   at (0,-8.2) {SCRIVI esito, scorte};
  \node[fc-term] (end)   at (0,-9.5) {FINE};
  \draw[fc-arrow] (start) -- (in);
  \draw[fc-arrow] (in) -- (dec);
  \draw[fc-arrow] (dec.west) -| node[fc-label, pos=0.25, above] {SÌ} (y1.north);
  \draw[fc-arrow] (dec.east) -| node[fc-label, pos=0.25, above] {NO} (n1.north);
  \draw[fc-arrow] (y1) -- (y2);
  \draw[fc-arrow] (y2.south) |- (join);
  \draw[fc-arrow] (n1.south) |- (join);
  \draw[fc-arrow] (join) -- (out);
  \draw[fc-arrow] (out) -- (end);
\end{tikzpicture}
\end{center}
La sottrazione delle scorte sta \textbf{solo} nel ramo SÌ: se il cibo non basta, il silos resta com'era.
\end{soluzione}

\begin{soluzione}[Esercizio 3: L'Indice di Felicità]
\begin{center}
\begin{tikzpicture}[font=\small]
  \node[fc-term] (start) at (0,0) {INIZIO};
  \node[fc-io]   (in)    at (0,-1.3) {LEGGI felicità};
  \node[fc-dec]  (dec1)  at (0,-2.9) {felicità $\geq$ 75 ?};
  \node[fc-op]   (y1)    at (-3.7,-4.6) {stato = ``FIORENTE \\ (+50 crediti)''};
  \node[fc-dec]  (dec2)  at (2.9,-4.6) {felicità $<$ 40 ?};
  \node[fc-op]   (y2)    at (0.6,-6.4) {stato = ``RIVOLTA''};
  \node[fc-op]   (n2)    at (5.2,-6.4) {stato = ``STABILE''};
  \node[fc-join] (join)  at (0,-7.8) {};
  \node[fc-io]   (out)   at (0,-8.9) {SCRIVI stato};
  \node[fc-term] (end)   at (0,-10.2) {FINE};
  \draw[fc-arrow] (start) -- (in);
  \draw[fc-arrow] (in) -- (dec1);
  \draw[fc-arrow] (dec1.west) -| node[fc-label, pos=0.25, above] {SÌ} (y1.north);
  \draw[fc-arrow] (dec1.east) -| node[fc-label, pos=0.25, above] {NO} (dec2.north);
  \draw[fc-arrow] (dec2.west) -| node[fc-label, pos=0.3, above] {SÌ} (y2.north);
  \draw[fc-arrow] (dec2.east) -| node[fc-label, pos=0.3, above] {NO} (n2.north);
  \draw[fc-arrow] (y1.south) |- (join);
  \draw[fc-arrow] (y2.south) |- (join);
  \draw[fc-arrow] (n2.south) |- (join);
  \draw[fc-arrow] (join) -- (out);
  \draw[fc-arrow] (out) -- (end);
\end{tikzpicture}
\end{center}
Il secondo rombo viene raggiunto solo quando la felicità è sotto 75: per questo il ramo NO finale copre esattamente l'intervallo da 40 a 74, senza bisogno di un terzo confronto.
\end{soluzione}

% ------------------------------------------------------------------------------
\section{Checklist del Flow Chart}
% ------------------------------------------------------------------------------

\begin{consiglio}[Prima di Consegnare il Disegno]
\begin{itemize}
  \item[$\square$] Ho usato i \textbf{parallelogrammi} per ogni \texttt{LEGGI} e \texttt{SCRIVI}?
  \item[$\square$] Ho usato i \textbf{rettangoli} per calcoli e assegnazioni?
  \item[$\square$] Ogni \textbf{rombo} ha due frecce uscenti con le etichette \textbf{SÌ} e \textbf{NO}?
  \item[$\square$] I rami di ogni bivio si \textbf{ricongiungono} prima dell'output finale?
  \item[$\square$] Il diagramma ha un solo \textbf{INIZIO} e una sola \textbf{FINE}?
\end{itemize}
\end{consiglio}
